%%% Курсовая работа — полный шаблон (все элементы)
%%% Тематика: Численные методы решения ОДУ
%%% Компиляция: xelatex → biber → xelatex → xelatex

\documentclass[font=modern]{cmcmsuthesis-cw}

\author{Иванов Иван Иванович}
\title{Численные методы решения обыкновенных
       дифференциальных уравнений}
\thesissetup{department = {Кафедра вычислительной математики}}
\supervisor{name = {канд.~физ.-мат.~наук, доцент Петров~П.П.}}

\usepackage{biblatex}
\addbibresource{refs.bib}

\begin{document}

\maketitle
\tableofcontents

% ================================================================
\section*{Введение}
\addcontentsline{toc}{section}{Введение}

Задача численного интегрирования ОДУ является фундаментальной
в вычислительной математике~\cite{shannon1948}.

\textbf{Целью} данной работы является сравнительный анализ методов
численного интегрирования ОДУ.

Для достижения цели необходимо решить следующие \textbf{задачи}:
\begin{enumerate}[beginpenalty=10000]
    \item реализовать методы Эйлера и Рунге--Кутты;
    \item исследовать устойчивость и точность;
    \item провести эксперименты на тестовых задачах;
    \item подготовить рекомендации по выбору метода.
\end{enumerate}

% ================================================================
\section{Теоретические основы}
\label{sec:theory}

\subsection{Основные определения}

Пусть $\vect{x}(t) \in \mathbb{R}^n$~--- вектор неизвестных.

\begin{definition}\label{def:cauchy}
\emph{Задачей Коши} называется задача нахождения решения
\begin{equation}\label{eq:ode}
    \frac{d\vect{x}}{dt} = \vect{f}(t, \vect{x}),
    \quad \vect{x}(t_0) = \vect{x}_0,
\end{equation}
где $\vect{f}\colon \mathbb{R} \times \mathbb{R}^n \to \mathbb{R}^n$.
\end{definition}

\begin{definition}\label{def:lipschitz}
Функция $\vect{f}$ удовлетворяет \emph{условию Липшица} по $\vect{x}$,
если $\exists L>0$: $\|\vect{f}(t,\vect{x})-\vect{f}(t,\vect{y})\|
\leq L\|\vect{x}-\vect{y}\|$.
\end{definition}

\begin{theorem}[Пикара--Линделёфа]\label{thm:picard}
Если $\vect{f}$ непрерывна и липшицева
(определение~\ref{def:lipschitz}), то задача
Коши~\eqref{eq:ode} имеет единственное решение.
\end{theorem}

\begin{proof}
Рассмотрим итерации Пикара
$\vect{x}^{(k+1)}(t) = \vect{x}_0 +
\int_{t_0}^{t}\vect{f}(s,\vect{x}^{(k)}(s))\,ds$
и покажем сходимость по лемме~\ref{lem:gronwall}.
\end{proof}

\begin{corollary}\label{cor:global}
Если $\vect{f}$ глобально липшицева, решение существует
на $[t_0,+\infty)$.
\end{corollary}

\begin{lemma}\label{lem:gronwall}
Если $u(t) \leq a + b\int_{t_0}^t u(s)\,ds$,
то $u(t) \leq a\,e^{b(t-t_0)}$.
\end{lemma}

\begin{proposition}\label{prop:euler-stable}
Метод Эйлера устойчив для $\dot{x}=\lambda x$
при $|1+h\lambda|<1$.
\end{proposition}

\begin{remark}\label{rem:stiff}
Для жёстких систем явные методы требуют малого шага $h$.
\end{remark}

\begin{example}\label{ex:harmonic}
Осциллятор $\ddot{x}+\omega^2 x=0$ сводится к системе:
$\dot{x}_1=x_2$, $\dot{x}_2=-\omega^2 x_1$.
\end{example}

\begin{problem}\label{prob:pendulum}
Исследовать нелинейный маятник $\ddot{\theta}+\sin\theta=0$.
\end{problem}

Команды пакета: $\matr{A}\in\mathbb{F}^{k\times n}_q$,
группа $\str{S}_n$, элемент $\elt{\sigma}$, код $\code{C}$,
алгоритм $\algo{RK4}$, сложность $\compl{O}(n^4)$,
шифр $\crypto{E}$.

% ================================================================
\section{Формулы}
\label{sec:formulas}

\subsection{Ненумерованные формулы}

Inline: $x \approx \sin x$ при $x \to 0$.
Display: \[(x_1+x_2)^2 = x_1^2 + 2x_1 x_2 + x_2^2.\]
Непрерывная дробь:
\[\frac{1}{\sqrt{2}+\displaystyle\frac{1}{\sqrt{2}+
    \displaystyle\frac{1}{\sqrt{2}+\cdots}}}\]

\subsection{Нумерованные формулы}
\begin{equation}\label{eq:euler-num}
e=\lim_{n\to\infty}\left(1+\frac{1}{n}\right)^{\!n}.
\end{equation}
\begin{equation}\label{eq:basel}
\lim_{n\to\infty}\sum_{k=1}^n\frac{1}{k^2}=\frac{\pi^2}{6}.
\end{equation}
Ссылки: \eqref{eq:euler-num} и~\eqref{eq:basel}.

С \texttt{multlined}:
\begin{equation}\label{eq:longsum}
\begin{multlined}
1+2+3+\dots+50+51+\dots+\\+96+97+98+99+100=5050.
\end{multlined}
\end{equation}

\subsection{Многострочные формулы}

\texttt{aligned}:
\[\begin{aligned}
f_W &= \min\!\left(1,\max\!\left(0,
    \frac{W/W_{\max}}{W_{\text{crit}}}\right)\right),\\
f_T &= \min\!\left(1,\max\!\left(0,
    \frac{T_s/T_{\text{melt}}}{T_{\text{crit}}}\right)\right).
\end{aligned}\]

\texttt{alignedat}:
\[|x|=\left\{\begin{alignedat}{2}
& & x,\quad &\text{если }x\geqslant 0;\\
&-& x,\quad &\text{если }x<0.
\end{alignedat}\right.\]

\texttt{cases}:
\[|x|=\begin{cases}\phantom{-}x,&\text{если }x\geqslant 0;\\
-x,&\text{если }x<0.\end{cases}\]

Нумерованная с \texttt{alignedat}:
\begin{equation}\label{eq:rk}
\begin{alignedat}{2}
\vect{k}_1 &= h\,\vect{f}(t_n,\vect{x}_n), &\quad&\\
\vect{k}_2 &= h\,\vect{f}(t_n+\tfrac{h}{2},
    \vect{x}_n+\tfrac{\vect{k}_1}{2}), &&\\
\vect{k}_3 &= h\,\vect{f}(t_n+\tfrac{h}{2},
    \vect{x}_n+\tfrac{\vect{k}_2}{2}), &&\\
\vect{k}_4 &= h\,\vect{f}(t_n+h,\vect{x}_n+\vect{k}_3).&&
\end{alignedat}
\end{equation}

\subsection{Субформулы}
\begin{subequations}\label{eq:sub1}
\begin{gather}
y = x^2+1 \label{eq:sub1a}\\
y = 2x^2-x+1 \label{eq:sub1b}\\
y = 3x^2-4x+8 \label{eq:sub1c}
\end{gather}
\end{subequations}
\begin{subequations}\label{eq:sub2}
\begin{align}
y &= x^3+x^2+x+1 \label{eq:sub2a}\\
y &= x^2 \label{eq:sub2b}
\end{align}
\end{subequations}
Ссылки: \eqref{eq:sub1a}, \eqref{eq:sub1b}, \eqref{eq:sub2a}, \eqref{eq:sub1c}.

Система уравнений (Лоренц):
\[\left\{\begin{array}{rl}
\dot x=&\sigma(y-x)\\\dot y=&x(r-z)-y\\\dot z=&xy-bz
\end{array}\right..\]

Матрица:
\[\left(\begin{array}{ccc}
a_{11}&\ldots&a_{1n}\\\vdots&\ddots&\vdots\\a_{n1}&\ldots&a_{nn}
\end{array}\right).\]

Длинная inline формула:
$a_1+a_2+a_3+a_4+a_5+a_6+a_7+a_8+a_9+a_{10}+a_{11}+a_{12}+a_{13}+a_{14}+a_{15}+a_{16}+a_{17}+a_{18}+a_{19}+a_{20}=\sum_{i=1}^{20}a_i.$

\subsection{Греческие буквы и алфавиты}

Обычные: $\alpha\beta\gamma\delta\epsilon\theta\lambda\mu\pi\sigma\phi\omega\Gamma\Delta\Omega$.

Прямые: $\symup{\alpha}\symup{\beta}\symup{\gamma}\symup{\delta}\symup{\Gamma}\symup{\Omega}$.

Жирные: $\symbf{\alpha}\symbf{\beta}\symbf{\gamma}\symbf{\Gamma}\symbf{\Omega}$.

Жирные прямые: $\symbfup{\alpha}\symbfup{\beta}\symbfup{\gamma}\symbfup{\Gamma}\symbfup{\Omega}$.

Алфавиты:
\[\begin{aligned}
\mathcal{ABCDEFGHIJKLMNOPQRSTUVWXYZ}\\
\mathfrak{ABCDEFGHIJKLMNOPQRSTUVWXYZ}\\
\mathbb{ABCDEFGHIJKLMNOPQRSTUVWXYZ}
\end{aligned}\]

Русские операторы: $\tan\alpha$/$\cot\beta$/$\csc\gamma$;
$\le$, $\leqslant$; $\phi$/$\varphi$; $\epsilon$/$\varepsilon$;
$\kappa$/$\varkappa$; $\emptyset$/$\varnothing$.

Отступы:
\[\begin{aligned}
\backslash! &\quad f(x)=x^2\!+3x\!+2\\
\text{default}&\quad f(x)=x^2+3x+2\\
\backslash, &\quad f(x)=x^2\,+3x\,+2\\
\backslash; &\quad f(x)=x^2\;+3x\;+2\\
\backslash\text{quad}&\quad f(x)=x^2\quad+3x\quad+2
\end{aligned}\]

% ================================================================
\section{Форматирование текста}

\begin{description}
\item[Жирный:] \textbf{жирный шрифт}.
\item[Курсив:] \emph{курсивный шрифт}.
\item[Жирный курсив:] \textbf{\emph{жирный курсив}}.
\item[Подчёркивание:] \underline{подчёркнутый текст}.
\item[Зачёркивание:] \sout{зачёркнутый текст}.
\item[Sans:] \textsf{текст без засечек, sans font}.
\end{description}

% ================================================================
\section{Таблицы}

\begin{table}[htbp]
\caption{Простая таблица}\label{tab:simple}
\centering
\begin{tabular}{|l|c|c|c|}
\hline Метод & $h$ & $h/2$ & $h/4$ \\\hline
Эйлер & 8.62 & 8.77 & 8.70 \\\hline
РК-2 & 7.90 & 7.93 & 7.93 \\\hline
РК-4 & 8.70 & 8.70 & 8.70 \\\hline
\end{tabular}
\end{table}

\begin{table}[htbp]
\caption{Таблица booktabs}\label{tab:booktabs}
\centering
\begin{tabular}{lrcc}
\toprule
\textbf{Метод}&\textbf{Погрешность}&\textbf{Время,с}&\textbf{Порядок}\\
\midrule
Эйлер & $1.2\times10^{-2}$ & 0.001 & 1\\
РК-4  & $5.6\times10^{-8}$ & 0.004 & 4\\
Адамс & $1.1\times10^{-6}$ & 0.003 & 4\\
\bottomrule
\end{tabular}
\end{table}

\begin{longtable}{@{}rr@{}}
\caption{Длинная таблица}\label{tab:long}\\
\toprule $t$&$x(t)$\\\midrule\endfirsthead
\toprule $t$&$x(t)$\\\midrule\endhead\bottomrule\endfoot
0.0&1.000\\0.1&1.105\\0.2&1.221\\0.3&1.350\\
0.4&1.492\\0.5&1.649\\0.6&1.822\\0.7&2.014\\
0.8&2.226\\0.9&2.460\\1.0&2.718\\
\end{longtable}

Ссылка на таблицу~\ref{tab:booktabs}.

% ================================================================
\section{Изображения}

%% Раскомментируйте при наличии файлов:
% \begin{figure}[htbp]\centering
% \includegraphics[width=0.8\textwidth]{images/plot.png}
% \caption{График решения}\label{fig:plot}
% \end{figure}
%
% \begin{figure}[htbp]\centering
% \includegraphics[scale=0.5]{images/logo.png}
% \caption{Логотип (масштаб 0.5)}\label{fig:logo}
% \end{figure}
%
% Ссылка на рисунок~\ref{fig:plot}.

% ================================================================
\section{Списки}

\subsection{Нумерованный}
\begin{enumerate}
\item Первый\begin{enumerate}
\item Второй\begin{enumerate}
\item Третий\begin{enumerate}
\item Четвёртый\end{enumerate}\end{enumerate}\end{enumerate}
\end{enumerate}

\subsection{Ненумерованный}
\begin{itemize}
\item Первый\begin{itemize}
\item Второй\begin{itemize}
\item Третий\begin{itemize}
\item Четвёртый\end{itemize}\end{itemize}\end{itemize}
\end{itemize}

% ================================================================
\section{Алгоритмы и листинги}

\begin{algorithm}[H]
\caption{Метод Рунге--Кутты}\label{alg:rk4}
\KwInput{$\vect{f}$, $t_0$, $\vect{x}_0$, $h$, $N$}
\KwOutput{массив решений}
\For{$i=1$ \KwTo $N$}{
    $\vect{k}_1\gets h\,\vect{f}(t,\vect{x})$\;
    $\vect{k}_2\gets h\,\vect{f}(t+h/2,\vect{x}+\vect{k}_1/2)$
        \tcp*{стадия 2}
    \Comment{обновляем}
    \uIf{$\|\vect{x}\|>10^{10}$}{\Return{ошибка}\;}
    \uElseIf{$h<10^{-15}$}{\Return{предупреждение}\;}
    \Else{сохранить $(t,\vect{x})$\;}
}
\Return{массив}\;
\end{algorithm}

Ссылка: алгоритм~\ref{alg:rk4}.

\begin{lstlisting}[language=Python,caption={Python},label={lst:py}]
import numpy as np
def euler(f, t0, x0, h, n):
    """Метод Эйлера для решения ОДУ."""
    t, x = t0, np.array(x0, dtype=float)
    for _ in range(n):
        x = x + h * f(t, x)  # Шаг Эйлера
        t += h
    return x
print("Результат:", euler(lambda t,x: -x, 0, [1.0], 0.1, 100))
\end{lstlisting}

\begin{lstlisting}[language=C,caption={C},label={lst:c}]
#include <stdio.h>
/* Правая часть ОДУ */
double f(double t, double x) { return -x; }
int main(void) {
    double t=0, x=1, h=0.01;
    // Интегрируем
    for (int i=0; i<1000; ++i) { x+=h*f(t,x); t+=h; }
    printf("x(%.1f) = %.6f\n", t, x);
    return 0;
}
\end{lstlisting}

\begin{lstlisting}[style=pseudocode,caption={Псевдокод}]
(*@\lstInput@*): $f$, $t_0$, $x_0$, $\varepsilon$
(*@\lstOutput@*): решение с точностью $\varepsilon$
procedure AdaptiveStep($f$, $t_0$, $x_0$, $\varepsilon$)
begin
    $h \gets 0.1$
    while $t < T$ do
    begin
        if $|x_1 - x_2| < \varepsilon$ then
            $t \gets t + h$
        else
            $h \gets h / 2$
    end
    return $x$
end
\end{lstlisting}

\begin{lstpseudocode}[caption={Бисекция},label={alg:bisect}]
(*@\lstInput@*): $[a,b]$, $g$, $\varepsilon$
(*@\lstOutput@*): корень $x^*$
function Bisection($a$, $b$, $g$, $\varepsilon$)
begin
    while $b - a > \varepsilon$ do
    begin
        $c \gets (a+b)/2$
        if $g(a)\cdot g(c)<0$ then $b\gets c$
        else $a\gets c$
    end
    return $(a+b)/2$
end
\end{lstpseudocode}

% ================================================================
\section{Библиография}

Одиночная~\cite{shannon1948}. Книга~\cite{knuth1984}.
Множественная~\cite{shannon1948,knuth1984,macwilliams1979}.

% ================================================================
\section*{Заключение}
\addcontentsline{toc}{section}{Заключение}

\begin{enumerate}
\item Реализованы методы Эйлера и Рунге--Кутты.
\item Сравнение приведено в таблице~\ref{tab:booktabs}.
\item Теоретические оценки подтверждены экспериментами.
\end{enumerate}

\clearpage\listoftables\addcontentsline{toc}{section}{Список таблиц}
\printbibliography

\appendix
\section{Полный код}
\subsection{Модуль интегрирования}
Код модуля.
\subsection{Модуль визуализации}
Код графиков.

\section{Дополнительные результаты}
Таблицы при разных шагах.
\subsection{Эксперимент 1}
Данные.

\end{document}
